\documentclass[12pt]{article}
\usepackage{amsmath, amssymb, amsthm}
\usepackage{mathtools}
\usepackage{tensor}
\usepackage{geometry}
\usepackage{booktabs}
\usepackage{array}
\usepackage{xcolor}
\usepackage{hyperref}
\usepackage{fancyhdr}
\usepackage{slashed}
\geometry{margin=1.1in, headheight=15pt}

\definecolor{cadblue}{RGB}{30,90,200}
\definecolor{verifygreen}{RGB}{20,140,60}

\newcommand{\cdbexpr}[1]{\textcolor{cadblue}{$#1$}}
\newcommand{\verified}[1]{\textcolor{verifygreen}{\checkmark\; #1}}
\newcommand{\half}{\tfrac{1}{2}}
\newcommand{\dal}{\dot{\alpha}}
\newcommand{\dbe}{\dot{\beta}}
\newcommand{\dga}{\dot{\gamma}}
\newcommand{\dde}{\dot{\delta}}
\newcommand{\sigmabar}{\bar{\sigma}}
\newcommand{\lam}{\lambda}
\newcommand{\lamtil}{\tilde{\lambda}}

\pagestyle{fancy}
\fancyhf{}
\lhead{Srednicki QFT --- Chapter 48}
\rhead{Massless Particles \& Spinor-Helicity}
\cfoot{\thepage}

\title{\textbf{Srednicki QFT: Chapter 48}\\[6pt]
       \large Massless Particles and Spinor-Helicity Formalism\\[4pt]
       \normalsize Cadabra2 expressions and numerical verification}
\author{Generated by \texttt{ch48\_export\_latex.py}}
\date{}

\begin{document}
\maketitle
\tableofcontents
\newpage

%% ============================================================
\section{§48.A\quad Massless Momentum as Spinor Outer Product}
%% ============================================================

\subsection{Rank-1 factorization}

A massless 4-momentum $p^\mu$ (with $p^2=0$) maps to a $2\times2$ spinor matrix
\begin{equation}
  p_{\alpha\dot\alpha} \;=\; p^\mu\,\sigma_{\mu,\,\alpha\dot\alpha},
  \qquad
  \sigma_\mu = (I_2,\,\boldsymbol{\sigma}),
  \tag{48.1}
\end{equation}
whose determinant equals $-p^2/4=0$ (for massless $p$), so the matrix has rank~1.
Any rank-1 $2\times2$ matrix factors as an outer product:
\begin{equation}
  \boxed{
    p_{\alpha\dot\alpha} = \lambda_\alpha\,\tilde\lambda_{\dot\alpha}
  }
  \qquad (p^2 = 0)
  \tag{48.2}
\end{equation}
The two-component \textbf{helicity spinors} $\lambda_\alpha$ (left-handed, undotted)
and $\tilde\lambda_{\dot\alpha}$ (right-handed, dotted) encode the full massless kinematics.

\medskip
\textbf{Cadabra2 representation:} \cdbexpr{\lambda_{\alpha} \tilde{\lambda}_{\dot{\alpha}}}

\subsection{Explicit components}

For $p^\mu = E(1,\,\sin\theta\cos\phi,\,\sin\theta\sin\phi,\,\cos\theta)$:
\begin{equation}
  \lambda_\alpha = \sqrt{2E}\begin{pmatrix}\cos\tfrac\theta2 \\ \sin\tfrac\theta2\,e^{i\phi}\end{pmatrix},
  \qquad
  \tilde\lambda_{\dot\alpha} = \sqrt{2E}\begin{pmatrix}\cos\tfrac\theta2 \\ \sin\tfrac\theta2\,e^{-i\phi}\end{pmatrix}
  \tag{48.3}
\end{equation}

\textbf{Numerical check} ($E=3$, $\theta=\pi/4$, $\phi=\pi/3$):
\[
  p_{\alpha\dot\alpha} = \begin{pmatrix}5.1213+0j & 1.0607-1.8371j \\ 1.0607+1.8371j & 0.8787+0j\end{pmatrix},
  \quad
  \lambda = (2.2630,\;0.4687+0.8118j),
  \quad
  \tilde\lambda = (2.2630,\;0.4687-0.8118j)
\]

\verified{$p_{\alpha\dot\alpha} = \lambda_\alpha\tilde\lambda_{\dot\alpha}$: max error $1.8\times10^{-15}$.}

\verified{$\det(p_{\alpha\dot\alpha}) = 5.7\times10^{-16}$ (should be 0, since $p^2=0$).}

%% ============================================================
\section{§48.B\quad Angle $\langle ij\rangle$ and Square $[ij]$ Brackets}
%% ============================================================

\subsection{Definitions}

The fundamental Lorentz-invariant bilinears are:
\begin{align}
  \langle i\,j\rangle &= \varepsilon^{\alpha\beta}\,\lambda^i_\alpha\,\lambda^j_\beta
  = \lambda^i_1\lambda^j_2 - \lambda^i_2\lambda^j_1
  & \text{(angle bracket, undotted)},
  \tag{48.4}\\[4pt]
  [i\,j] &= \varepsilon_{\dot\alpha\dot\beta}\,\tilde\lambda^{i\,\dot\alpha}\tilde\lambda^{j\,\dot\beta}
  = \tilde\lambda^i_{\dot1}\tilde\lambda^j_{\dot2} - \tilde\lambda^i_{\dot2}\tilde\lambda^j_{\dot1}
  & \text{(square bracket, dotted)},
  \tag{48.5}
\end{align}
where $\varepsilon^{12}=+1$ (Srednicki convention).

\textbf{Cadabra2:}\quad
\cdbexpr{\abra{k}{p}} $= \langle k\,p\rangle$\quad and\quad
\cdbexpr{\sbra{k}{p}} $= [k\,p]$

\subsection{Antisymmetry}

Both brackets are antisymmetric in their arguments:
\begin{equation}
  \langle i\,j\rangle = -\langle j\,i\rangle,
  \qquad
  [i\,j] = -[j\,i].
  \tag{48.6}
\end{equation}

\verified{$|\langle kq\rangle + \langle qk\rangle| = 0.0\times10^0$ (exact, by construction).}
\verified{$|[kq] + [qk]| = 0.0\times10^0$ (exact, by construction).}

\subsection{Schouten identity}

Because spinor space is 2-dimensional, any three undotted spinors are linearly
dependent.  This yields the \textbf{Schouten identity}:
\begin{equation}
  \boxed{
    \langle i\,j\rangle\langle k\,l\rangle
    + \langle i\,k\rangle\langle l\,j\rangle
    + \langle i\,l\rangle\langle j\,k\rangle = 0
  }
  \tag{48.7}
\end{equation}
and identically for square brackets.

\textbf{Cadabra2 (LHS sum):} \cdbexpr{\abra{i}{j}\abra{k}{l} + \abra{i}{k}\abra{l}{j} + \abra{i}{l}\abra{j}{k}}

\verified{$|\langle12\rangle\langle34\rangle + \langle13\rangle\langle42\rangle + \langle14\rangle\langle23\rangle| = 4.6\times10^{-16}$ (four generic massless momenta).}

%% ============================================================
\section{§48.C\quad Mandelstam Variables from Spinors}
%% ============================================================

For massless momenta $p_i$, $p_j$ (with $p_i^2=p_j^2=0$), the kinematic invariant is:
\begin{equation}
  \boxed{
    s_{ij} = (p_i+p_j)^2 = 2\,p_i\cdot p_j = \langle i\,j\rangle[j\,i]
  }
  \tag{48.8}
\end{equation}

\textbf{Cadabra2:} \cdbexpr{\abra{k}{p}\sbra{p}{k}} $= s_{kp}$

\medskip
\textbf{Physical meaning:} The product $\langle kq\rangle[qk]$ gives the
Mandelstam invariant directly from the helicity spinors, without needing to
evaluate $p^\mu q_\mu$ with the metric.

\textbf{Numerical check} ($k=(2.5,0,0,2.5)$, $q=3(1,\sin(\pi/3),0,\cos(\pi/3))$):
\begin{align*}
  \langle kq\rangle[qk] &= 7.500000 \\
  -2\,k\cdot q &= 7.500000
\end{align*}

\verified{$|\langle kq\rangle[qk] - (-2k\cdot q)| = 8.9\times10^{-16}$.}

%% ============================================================
\section{§48.D\quad Helicity Spinors $u_\pm(k)$ in the Chiral Basis}
%% ============================================================

\subsection{Z-axis momentum: explicit form}

For massless momentum $k^\mu=(\omega,0,0,\omega)$ along the $+z$ axis,
the massless Dirac equation $\slashed{k}|\psi\rangle=0$ has solutions
(Srednicki eq.~60.10):
\begin{equation}
  |k] = u_-(k) = \sqrt{2\omega}\,(0,1,0,0)^T,
  \qquad
  |k\rangle = u_+(k) = \sqrt{2\omega}\,(0,0,1,0)^T
  \tag{48.9}
\end{equation}
where $|k]$ has negative helicity ($h=-\tfrac12$) and $|k\rangle$ has positive helicity ($h=+\tfrac12$).

The corresponding 2-component helicity spinors:
\begin{equation}
  \lambda_\alpha(k)\big|_{z\text{-axis}} = \sqrt{2\omega}\begin{pmatrix}1\\0\end{pmatrix},
  \qquad
  \tilde\lambda_{\dot\alpha}(k)\big|_{z\text{-axis}} = \sqrt{2\omega}\begin{pmatrix}1\\0\end{pmatrix}
  \tag{48.10}
\end{equation}

\verified{Massless Dirac eq.\ $\slashed{k}|k] = 0$: max error $0.0\times10^0$ (exact).}
\verified{Massless Dirac eq.\ $\slashed{k}|k\rangle = 0$: max error $0.0\times10^0$ (exact).}

\subsection{General massless momentum}

For $p^\mu = E(1,\sin\theta\cos\phi,\sin\theta\sin\phi,\cos\theta)$, the
4-component spinors in the Weyl representation are:
\begin{equation}
  |p] = \sqrt{2E}\begin{pmatrix}-\sin\tfrac\theta2\,e^{-i\phi} \\ \cos\tfrac\theta2 \\ 0 \\ 0\end{pmatrix},
  \qquad
  |p\rangle = \sqrt{2E}\begin{pmatrix}0 \\ 0 \\ \cos\tfrac\theta2 \\ \sin\tfrac\theta2\,e^{i\phi}\end{pmatrix}
  \tag{48.11}
\end{equation}
These satisfy $\slashed{p}|p]=\slashed{p}|p\rangle=0$ and are eigenstates of
the helicity operator $\hat h = \hat{\mathbf{J}}\cdot\hat{\mathbf{p}}$.

%% ============================================================
\section{§48.E\quad Gluon Polarization Vectors as Spinor Bilinears}
%% ============================================================

\subsection{Definitions}

The photon/gluon polarization vectors in the spinor-helicity formalism
(Srednicki between eqs.~60.6--60.9):
\begin{equation}
  \boxed{
    \varepsilon_+^\mu(k;\,q) = \frac{\tilde\lambda_q^{\dagger\dot\alpha}\,\sigma^\mu_{\alpha\dot\alpha}\,\lambda_k^\alpha}{\sqrt{2}\,\langle q\,k\rangle},
    \qquad
    \varepsilon_-^\mu(k;\,q) = \bigl[\varepsilon_+^\mu(k;\,q)\bigr]^*
  }
  \tag{48.12}
\end{equation}
Here $k$ is the photon momentum and $q\neq k$ is an arbitrary massless
\textbf{reference momentum} encoding residual gauge freedom.

\textbf{Ward identity (transversality):}
\begin{equation}
  k_\mu\,\varepsilon_\pm^\mu(k;\,q) = 0
  \tag{48.13}
\end{equation}

\textbf{Normalization:}
\begin{equation}
  \varepsilon_+^\mu\varepsilon_{+\mu} = \varepsilon_-^\mu\varepsilon_{-\mu} = -1,
  \qquad
  \varepsilon_+^\mu\varepsilon_{-\mu} = 0
  \tag{48.14}
\end{equation}

\subsection{Numerical verification}

Setup: $k^\mu=(2,0,0,2)$ along $+z$;
reference $q$ at $\theta=\pi/3$, $\phi=\pi/4$ with $\omega_q=1.5$.

Computed: $\varepsilon_+^\mu(k;q) \approx (-0.866+0.866i,\,-0.707,\,-0.707i,\,-0.866+0.866i)$

\verified{$|k_\mu\varepsilon_+^\mu(k;q)| = 0.0\times10^0$ (Ward identity for $\varepsilon_+$).}
\verified{$|k_\mu\varepsilon_-^\mu(k;q)| = 0.0\times10^0$ (Ward identity for $\varepsilon_-$).}

%% ============================================================
\section{§48.F\quad Little Group Scaling and Helicity Weights}
%% ============================================================

\subsection{The little group}

The little group of a massless momentum $p^\mu$ is $U(1)\simeq SO(2)$:
rotations about the direction of motion.  Under the little-group transformation
\begin{equation}
  \lambda_\alpha \;\to\; t\,\lambda_\alpha,
  \qquad
  \tilde\lambda_{\dot\alpha} \;\to\; t^{-1}\,\tilde\lambda_{\dot\alpha},
  \qquad t\in\mathbb{C}^*,
  \tag{48.15}
\end{equation}
the momentum $p_{\alpha\dot\alpha}=\lambda_\alpha\tilde\lambda_{\dot\alpha}$ is
invariant.  A massless particle state of helicity $h$ transforms as:
\begin{equation}
  |p,h\rangle \;\to\; t^{-2h}\,|p,h\rangle
  \tag{48.16}
\end{equation}

The spinor brackets transform with definite weights under particle $i$'s little group ($t_i=t$):
\begin{equation}
  \langle i\,j\rangle \;\to\; t_i\,t_j\,\langle i\,j\rangle,
  \qquad
  [i\,j] \;\to\; t_i^{-1}\,t_j^{-1}\,[i\,j].
  \tag{48.17}
\end{equation}
Under $\lambda_k\to t\lambda_k$ (with $\tilde\lambda_k\to t^{-1}\tilde\lambda_k$ to
preserve $p$), the polarization vector $\varepsilon_+^\mu(k;q)$ is invariant as a
4-vector.  The helicity weight manifests in the full amplitude
$\mathcal{A} \to \prod_i t_i^{-2h_i}\,\mathcal{A}$ (on-shell condition).

\verified{$\varepsilon_+^\mu(k;q)$ invariant under $\lambda_k\to t\lambda_k$,
$\tilde\lambda_k\to t^{-1}\tilde\lambda_k$ (preserving $p_k$): ratio $= 1.0000$.}

%% ============================================================
\section{§48.G\quad MHV Parke-Taylor Formula (Seed)}
%% ============================================================

\subsection{Motivation}

The spinor-helicity formalism makes tree-level gluon amplitudes remarkably
compact.  Amplitudes with all-same or all-but-one helicities vanish.
The first nonzero class are the \textbf{MHV (Maximally Helicity Violating)}
amplitudes, with exactly two negative helicities.

\subsection{Parke-Taylor formula}

The $n$-gluon colour-ordered MHV amplitude (two negative helicities at positions $i,j$):
\begin{equation}
  \boxed{
    \mathcal{A}_n^{\rm MHV}\!\bigl(1^+,\ldots,i^-,\ldots,j^-,\ldots,n^+\bigr)
    = \frac{\langle i\,j\rangle^4}
           {\langle 1\,2\rangle\langle 2\,3\rangle\cdots\langle n\,1\rangle}
  }
  \tag{48.18}
\end{equation}
(overall coupling and colour factors suppressed).
This formula, the \textbf{Parke-Taylor formula}, is the \emph{seed} for BCFW
on-shell recursion: all tree amplitudes can be built from it by shifting two
complex spinors and summing over factorisation channels.

\subsection{Structure check (3-point MHV)}

The $n=3$ MHV amplitude
$\mathcal{A}_3(1^-,2^-,3^+) = \langle12\rangle^3/(\langle12\rangle\langle23\rangle\langle31\rangle)$
is a pure ratio of angle brackets.  Numerically with generic momenta
$(E=1,\,\theta_1=0,\,\theta_2=\pi/3,\,\theta_3=2\pi/3)$:
\[
  \mathcal{A}_3 = \frac{\langle12\rangle^3}{\langle12\rangle\langle23\rangle\langle31\rangle}
  = \frac{(0.7071+0.7071i)^3}{(0.7071+0.7071i)(-0.3536+1.1464i)(-0.0000-1.7321i)}
  = 0.1418+0.4599i
\]

\verified{Parke-Taylor $n=3$ ratio evaluated from spinors (dimensionless Lorentz invariant).}

%% ============================================================
\section{Numerical Verification Summary}
%% ============================================================

\begin{center}
\renewcommand{\arraystretch}{1.5}
\begin{tabular}{@{}lll@{}}
  \toprule
  Identity & Equation & Max error \\
  \midrule
  $p_{\alpha\dot\alpha} = \lambda_\alpha\tilde\lambda_{\dot\alpha}$ (rank-1)
    & (48.2) & $1.8\times10^{-15}$ \\
  $\det(p_{\alpha\dot\alpha}) = 0$
    & ($p^2=0$) & $5.7\times10^{-16}$ \\
  $\langle ij\rangle = -\langle ji\rangle$ (antisymmetry)
    & (48.6) & $0$ (exact) \\
  $[ij] = -[ji]$ (antisymmetry)
    & (48.6) & $0$ (exact) \\
  Schouten identity $\langle\cdot\rangle$ brackets
    & (48.7) & $4.6\times10^{-16}$ \\
  $\langle kq\rangle[qk] = -2k\cdot q$
    & (48.8) & $8.9\times10^{-16}$ \\
  $\slashed{k}|k] = 0$ (massless Dirac)
    & (48.9) & $0$ (exact) \\
  $\slashed{k}|k\rangle = 0$ (massless Dirac)
    & (48.9) & $0$ (exact) \\
  $k_\mu\varepsilon_+^\mu = 0$ (Ward)
    & (48.13) & $0$ (exact) \\
  $k_\mu\varepsilon_-^\mu = 0$ (Ward)
    & (48.13) & $0$ (exact) \\
  $\varepsilon_+^\mu$ little-group invariant ($\lambda_k\to t\lambda_k$)
    & (48.15--16) & $< 10^{-15}$ \\
  \bottomrule
\end{tabular}
\end{center}

All identities verified to machine precision ($\lesssim 10^{-15}$).

%% ============================================================
\section{Summary}
%% ============================================================

\begin{center}
\renewcommand{\arraystretch}{1.5}
\begin{tabular}{@{}ll@{}}
  \toprule
  Object & Expression \\
  \midrule
  Spinor momentum & $p_{\alpha\dot\alpha} = \lambda_\alpha\tilde\lambda_{\dot\alpha}$;\quad $p^2=0$ iff rank 1 \\
  Angle bracket & $\langle ij\rangle = \varepsilon^{\alpha\beta}\lambda^i_\alpha\lambda^j_\beta = -\langle ji\rangle$ \\
  Square bracket & $[ij] = \varepsilon_{\dot\alpha\dot\beta}\tilde\lambda^{i\dot\alpha}\tilde\lambda^{j\dot\beta} = -[ji]$ \\
  Mandelstam & $s_{ij} = \langle ij\rangle[ji] = -2p_i\cdot p_j$ \\
  $z$-axis spinors & $|k]=\sqrt{2\omega}(0,1,0,0)^T$;\quad $|k\rangle=\sqrt{2\omega}(0,0,1,0)^T$ \\
  Polarization $\varepsilon_+$ & $\varepsilon_+^\mu(k;q) = \tilde\lambda_q^{\dagger}\sigma^\mu\lambda_k / (\sqrt{2}\langle qk\rangle)$ \\
  Polarization $\varepsilon_-$ & $\varepsilon_-^\mu(k;q) = [\varepsilon_+^\mu(k;q)]^*$ \\
  Ward identity & $k_\mu\varepsilon_\pm^\mu = 0$ \\
  Little group & $\lambda\to t\lambda$,\; $\tilde\lambda\to t^{-1}\tilde\lambda$;\; amplitude $\to t^{-2h}$ \\
  Schouten & $\langle ij\rangle\langle kl\rangle+\langle ik\rangle\langle lj\rangle+\langle il\rangle\langle jk\rangle=0$ \\
  MHV (Parke-Taylor) & $\mathcal{A}_n^{\rm MHV}=\langle ij\rangle^4/(\langle12\rangle\cdots\langle n1\rangle)$ \\
  \bottomrule
\end{tabular}
\end{center}

\bigskip
\noindent\textbf{Next:} Chapter 60 applies these tools to spinor electrodynamics:
twistor notation $|p], |p\rangle, [p|, \langle p|$, polarization vectors as
spinor bilinears, and the Fierz identities for loop computations.

\end{document}
